\clearpage\thispagestyle{empty}
\pdfbookmark[0]{Quellen zur Florabeobachtung vom 5. Januar 1848}{flora625-sources}
\begin{center}\textbf{Quellen zur Florabeobachtung vom 5. Januar 1848}\end{center}
\noindent Diese KI-Übertragung ins heutige Deutsch folgt dem Auszug auf S.~476 der \href{https://gdz.sub.uni-goettingen.de/id/PPN236007467}{\textit{Werke}, Band~VI} (Göttingen, 1874). Carl Friedrich Gauß schrieb am 10. Januar 1848 an Heinrich Christian Schumacher über die Beobachtung der Flora vom 5. Januar.
\par\smallskip\noindent Der Auszug erschien in den \textit{Astronomischen Nachrichten}, Band~XXVII, Nr.~625, \href{https://api.digitale-sammlungen.de/iiif/image/v2/bsb10049452_00015/full/1800,/0/default.jpg}{Sp.~13--14} und \href{https://api.digitale-sammlungen.de/iiif/image/v2/bsb10049452_00016/full/1800,/0/default.jpg}{Sp.~15--16}, ausgegeben am 29. Januar 1848. Gauß vergleicht die Beobachtungsbedingungen am 3. und 5. Januar und gibt zwei Abweichungen für Goldschmidts Elemente an. KI-Übertragung und KI-Satz.
\clearpage
